% SPDX-License-Identifier: GPL-3.0-or-later
% Copyright (C) 2026 Alessandro Greco
% Общая разметка и стили TikZ; внутренние имена окружений сохраняются.
% SPDX-License-Identifier: GPL-3.0-or-later
% Copyright (C) 2026 Alessandro Greco
% Colori
\definecolor{torviola}{HTML}{7D4698}
\definecolor{torverde}{HTML}{68B030}
\definecolor{torgrigio}{HTML}{4A4A4A}
\definecolor{avviso}{HTML}{C0392B}
\definecolor{sfondocodice}{HTML}{F6F4F8}

% Comandi di markup semantico
\newcommand{\cella}[1]{\texttt{\textbf{#1}}\index{#1@\texttt{#1}}}
\newcommand{\opt}[1]{\texttt{#1}\index{torrc!#1@\texttt{#1}}}
% \file: nomi di file e percorsi in monospaziato; possono andare a capo dopo "_" e dopo "/"
% (i percorsi di labs/it/ sono lunghi). Comando robusto: funziona anche nelle didascalie.
\ExplSyntaxOn
\tl_new:N \l__libro_file_tl
\NewDocumentCommand{\file}{m}
  {
    \group_begin:
    \tl_set:Nn \l__libro_file_tl {#1}
    \tl_replace_all:Nnn \l__libro_file_tl { / } { / \allowbreak }
    \cs_set_eq:NN \__libro_trattino_basso: \_
    \cs_set:Npn \_ { \__libro_trattino_basso: \allowbreak }
    \texttt { \tl_use:N \l__libro_file_tl }
    \group_end:
  }
\ExplSyntaxOff
\newcommand{\cmd}[1]{\texttt{#1}}
% URL del repository del libro (codice, laboratori, rete di prova).
\newcommand{\repolibro}{\url{https://github.com/aleff-github/onion-routing-anatomy}}
% Titolo di un'opera citata (articolo, proposta, libro): escluso dal controllo di originalità.
\newcommand{\titolo}[1]{\emph{#1}}
\newcommand{\termine}[2]{\emph{#1} (#2)\index{#1}}
\newcommand{\spec}[2]{\textsf{[#1~§\,#2]}}
\newcommand{\proposta}[1]{Proposta~#1\index{Proposte!#1}}
\newcommand{\aggiornatoal}[2]{%
  \begin{center}\small\color{torgrigio}\textit{Aggiornato al #1 --- #2}\end{center}}
\newcommand{\dacontrollare}[1]{\textcolor{avviso}{\textbf{[DA CONTROLLARE: #1]}}}
% Testo alternativo di una figura (accessibilita'). Nella stampa non produce
% nulla; nella conversione EPUB (tex4ht) viene ridefinito in libro/epub.tex per
% assegnare l'attributo alt dell'immagine. Va messo prima dell'ambiente grafico.
\newcommand{\altfig}[1]{}

% Riquadri
\tcbset{base/.style={enhanced,breakable,boxrule=0.6pt,arc=2pt,left=6pt,right=6pt,fonttitle=\bfseries\sffamily}}
\newtcolorbox{obiettivi}{base,colback=torviola!5,colframe=torviola,title={Obiettivi del capitolo}}
\newtcolorbox{approfondimento}[1][]{base,colback=torgrigio!4,colframe=torgrigio,title={Approfondimento\ifx&#1&\else: #1\fi}}
\newtcolorbox{attenzione}[1][]{base,colback=avviso!5,colframe=avviso,title={Attenzione\ifx&#1&\else: #1\fi}}
\newtcolorbox{storia}[1][]{base,colback=torverde!5,colframe=torverde!70!black,title={Nota storica\ifx&#1&\else: #1\fi}}
\newtcolorbox{specbox}[1]{base,colback=white,colframe=torviola!60,title={Dalla specifica: #1}}

% Laboratori ed esercizi
% Laboratori: [N] per i laboratori numerati di torlab (Parte II), altrimenti numerazione per capitolo.
\newcounter{laboratorio}[chapter]
\newtcolorbox{laboratorio}[2][]{base,colback=torverde!4,colframe=torverde!60!black,
  title={\ifx&#1&\refstepcounter{laboratorio}Laboratorio \thechapter.\arabic{laboratorio}\else Laboratorio~#1\fi: #2}}
\newcounter{esercizio}[chapter]
\renewcommand{\theesercizio}{\thechapter.\arabic{esercizio}}
\newcommand{\stelle}[1]{\ifcase#1\or$\star$\or$\star\star$\or$\star\star\star$\fi}
\newenvironment{esercizio}[1][1]{\refstepcounter{esercizio}\par\medskip
  \noindent\textbf{Esercizio \theesercizio}\ \stelle{#1}\quad}{\par}
\newenvironment{soluzione}[1]{\par\medskip\noindent\textbf{Soluzione \ref{#1}.}\quad}{\par}

% Codice
\lstset{basicstyle=\ttfamily\small,backgroundcolor=\color{sfondocodice},frame=single,
  rulecolor=\color{torviola!40},breaklines=true,columns=fullflexible,keepspaces=true,
  % senza questa, listings stampa gli spazi dentro le stringhe come \textvisiblespace:
  % il codice diventa illeggibile e non copiabile.
  showstringspaces=false,
  numbers=left,numberstyle=\tiny\color{torgrigio},keywordstyle=\color{torviola}\bfseries,
  commentstyle=\color{torgrigio}\itshape,stringstyle=\color{torverde!60!black},
  literate={à}{{\`a}}1 {è}{{\`e}}1 {é}{{\'e}}1 {ì}{{\`i}}1 {ò}{{\`o}}1 {ù}{{\`u}}1}
% listings conosce i caratteri solo fino a U+00FF: con LuaTeX il cirillico e il cinese
% dei listati uscivano fuori posto e perdevano gli spazi, nel PDF e nell'EPUB. Mentre si
% legge un listato, ../shared/lst-unicode.lua racchiude quei caratteri fra due ` (il
% carattere di escape, assente da tutti i listati) e li fa comporre come testo normale.
\lstset{escapechar=`}
\directlua{lstunicode = dofile("../shared/lst-unicode.lua")}
\AddToHook{env/lstlisting/begin}{\directlua{lstunicode.enable()}}
\AddToHook{env/lstlisting/end}{\directlua{lstunicode.disable()}}
\lstdefinelanguage{torrc}{morekeywords={SocksPort,ControlPort,ORPort,DirPort,ExitPolicy,
  HiddenServiceDir,HiddenServicePort,UseBridges,Bridge,ClientTransportPlugin,Nickname,ContactInfo},
  morecomment=[l]{\#}}

% Stili TikZ comuni
\tikzset{
  client/.style={draw=torviola,fill=torviola!10,rounded corners,minimum width=1.8cm,minimum height=0.9cm,font=\sffamily\small},
  relay/.style={draw=torverde!60!black,fill=torverde!12,circle,minimum size=1.1cm,font=\sffamily\small},
  server/.style={draw=torgrigio,fill=torgrigio!8,minimum width=1.8cm,minimum height=0.9cm,font=\sffamily\small},
  cellaf/.style={-{Stealth[length=2.5mm]},thick,torviola},
  risposta/.style={-{Stealth[length=2.5mm]},thick,dashed,torverde!60!black},
}

% Необязательный аргумент — падежная форма: \proposta[Предложения]{289} («из Предложения 289»).
\renewcommand{\proposta}[2][Предложение]{#1~#2\index{Предложения Tor!#2}}
% Сохраняем исходный английский термин из аргумента #2 для читателя.
\renewcommand{\termine}[2]{\emph{#1} (#2)\index{#1}}
\renewcommand{\aggiornatoal}[2]{%
  \begin{center}\small\color{torgrigio}\textit{Последняя проверка: #1 --- #2}\end{center}}
\renewcommand{\dacontrollare}[1]{\textcolor{avviso}{\textbf{[ПРОВЕРИТЬ: #1]}}}
\renewtcolorbox{obiettivi}{base,colback=torviola!5,colframe=torviola,title={Цели главы}}
\renewtcolorbox{approfondimento}[1][]{base,colback=torgrigio!4,colframe=torgrigio,title={Подробнее\ifx&#1&\else: #1\fi}}
\renewtcolorbox{attenzione}[1][]{base,colback=avviso!5,colframe=avviso,title={Внимание\ifx&#1&\else: #1\fi}}
\renewtcolorbox{storia}[1][]{base,colback=torverde!5,colframe=torverde!70!black,title={Историческая справка\ifx&#1&\else: #1\fi}}
\renewtcolorbox{specbox}[1]{base,colback=white,colframe=torviola!60,title={Из спецификации: #1}}
\renewtcolorbox{laboratorio}[2][]{base,colback=torverde!4,colframe=torverde!60!black,
  title={\ifx&#1&\refstepcounter{laboratorio}Лабораторная работа \thechapter.\arabic{laboratorio}\else Лабораторная работа~#1\fi: #2}}
\renewenvironment{esercizio}[1][1]{\refstepcounter{esercizio}\par\medskip
  \noindent\textbf{Упражнение \theesercizio}\ \stelle{#1}\quad}{\par}
\renewenvironment{soluzione}[1]{\par\medskip\noindent\textbf{Решение \ref{#1}.}\quad}{\par}
